% Network methods used inside the manuscript's \matmethods{...} block.
\subsection*{Ingredient network and recipe scores}
We lowercased ingredient names and merged 72 manually reviewed aliases. We excluded 20 generic names, including salt, pepper, water, and unspecified oil, and names found in fewer than three or more than 35\% of recipes. The resulting binary recipe-by-ingredient matrix has 3,877 ingredient types. An entry of one is a link in the bipartite graph; the ingredient projection joins two ingredients when they appear together in at least one recipe. To describe ingredient concentration across cuisine tags, we calculated each ingredient's entropy over the nine tags, divided by $\log 9$, and averaged across ingredients with their recipe counts as weights. We compared this with 199 shuffles of the recipe cuisine labels on the fixed bipartite network.

For cuisine $c$, let $N_c$ be its number of recipes and $n_{c,ij}$ the number containing both ingredients $i$ and $j$. Following Ahn et al. \cite{ahn2011flavor}, we define the pair contrast as $n_{c,ij}/N_c - \langle n_{d,ij}/N_d\rangle_{d\ne c}$, an equal-weight mean over the other eight cuisines. The pair score for a recipe is the mean contrast over all its retained pairs. To prevent a recipe from increasing its own score, we replace its cuisine term with $(n_{c,ij}-1)/(N_c-1)$ when scoring it. We calculate the ingredient-only score in the same way, replacing pairs with individual ingredients. Recipes with fewer than two retained ingredients have no pair score. For illustrative pair signatures, we require at least 20 recipes containing the pair in the named cuisine; this threshold does not affect recipe scores. We exported both scores with recipe IDs for the later review analysis.

\subsection*{Pairing score}
Within cuisine $c$, let $n_{c,ij}$ be the number of other recipes containing both $i$ and $j$, and $E_{c,ij}$ its mean over 50 Curveball randomizations of the recipe--ingredient matrix within each cuisine (ten attempted trades per recipe), which preserve every recipe's ingredient count and every ingredient's frequency in the cuisine. For leave-one-out, $E_{c,ij}$ is rescaled by $(n_{c,i}-1)(n_{c,j}-1)/(n_{c,i}n_{c,j})$. The pairing lift is $L_{c,ij}=\log_2\left[(n_{c,ij}+k)/(E_{c,ij}+k)\right]$, positive when two ingredients are combined more often than their popularity predicts, and a recipe's pairing score is the mean $L_{c,ij}$ over its retained pairs. This is the network analogue of a collocation score in text. The pseudocount $k$ shrinks rare pairs, for which a count of one or two is mostly chance: 34\% of all recipe pairs have $E_{c,ij}<1$. We chose $k=5$ by split-half reliability: we split each cuisine's recipes at random into two halves, scored every recipe once from each half, and correlated the two scores. The correlation rose from 0.70 with $k=1$ to 0.82 with $k=5$ (Spearman--Brown reliability with all data: 0.82 to 0.90). Scores with $k=2$ and $k=10$ correlate at $r\geq0.98$ with $k=5$, and we report them as robustness checks.

\subsection*{Supporting network comparisons}
We also projected the bipartite graph onto recipes. We measured similarity with $J(r,s)=|I_r\cap I_s|/|I_r\cup I_s|$ and linked recipes sharing at least two retained ingredients with $J\geq0.18$. We kept the union of each recipe's 15 most similar neighbors; link weight was $J$. Cuisine labels played no role in constructing links. Weighted Louvain community detection was run without cuisine labels; we inspected ingredient and title patterns in its largest mixed community. To compare the observed within-cuisine weight share with cuisine sizes alone, we shuffled cuisine labels 199 times over the fixed graph. We then ran 49 Curveball randomizations of the recipe--ingredient matrix within each cuisine, with ten attempted trades per recipe, preserving every recipe's ingredient count and each ingredient's frequency within its cuisine. Another 49 runs restricted trades to recipes of the same cuisine and broad Food.com course type. We rebuilt the recipe projection after every randomization and checked the preserved counts. The first shuffle tests an association between tags and the fixed graph; the ingredient swaps ask how much structure remains after accounting for ingredient frequencies and broad dish types.

For the nine illustrative pairs, we used a separate set of 49 cuisine-by-course-type Curveball randomizations with ten attempted trades per recipe. We counted how many recipes in the named cuisine contained both ingredients in each randomized matrix. We checked that row sizes and ingredient frequencies in every cuisine-by-type group were unchanged after every run. The middle 95\% of randomized counts is a reference range for these nine previously selected pairs, not a confidence interval or a multiple-comparison test.
